\scanpage{209}
rôle de \uline{classes de chemins}, ils sont déterminés \inserted{(s'il y en a au moins un)} \unsure{mod} automorphismes \unsure{du premier|des premiers}, qui font les p\textsuperscript{ts} d'un groupe algébrique. Mais en général il peut y avoir des p\textsuperscript{ts} à valeurs dans un $=$ $k$ non isomorphes. Les groupes de Galois \unsure{opèrent|opérant|opère} sur \deleted{les points} \unsure{l'en\textsuperscript{s}|l'un|l'ens.} des p\textsuperscript{ts} à valeurs dans $k$. Si on a un pt $\unsure{\xi|\zeta}$ à val \unsure{dans} $k$, tous les autres p\textsuperscript{ts} s'en déduisent par des fibrés principaux \unsure{hom.} sur le (pro)groupe algébrique $\underline{\mathrm{Aut}}(\xi)$.

\diagram{\[
\begin{array}{@{}l@{}}
\left\{\begin{array}{ll}
\text{Pts $\ell$-adiques} & \text{valeurs }\mathbb Z_{\ell}\text{ ou }\mathbb Q_{\ell}\\
\text{pts transcendants} & \text{valeurs }\mathbb R\ \deleted{\text{ou }\mathbb C}\ \deleted{\unsure{\text{déjà}}}\text{, voire }\mathbb Z\\
 & \text{[structure supplémentaire]}\\
 & \text{Il semble que ces p\textsuperscript{ts} correspondent}\\
 & \text{\unsure{biunivoq\textsuperscript{mt}|biunivoquent|biunivoquement} aux pts de}\\
 & X\text{ à valeurs dans }\mathbb C\text{,}\\
 & \text{mod conjugaison}
\end{array}\right.\\[0.5em]
\begin{array}{ll}
\text{pts Hodge-de Rham} & \text{\unsure{valeur|valeurs} }k\text{ \unsure{quelconques|quelconque}}\\
\text{pts de Witt} & \text{valeurs }W_{p}(k)\ \inserted{\text{$k$ corps de car $p$}}
\end{array}
\end{array}
\]}
% archive sheet number: 208

